\section{Tangent sheaves and Hermitian--Einstein limits}\label{sec:tangent}

The tangent bundles provide compatible nontrivial holomorphic data on the
same transition. Their Hermitian--Einstein connections are intrinsic to
the Ricci-flat metrics, but their algebraic extensions across the central
fibre require a separate calculation. We make that distinction explicit.

\subsection{The flat limit and its reflexive hull}
For the family \eqref{eq:split-equations}, put
\[
 \mathcal A=T_{\PP^4}|_{\mathcal X},\qquad
 \mathcal B=\OO_{\mathcal X}(5),\qquad
 \mathcal K=\ker(\mathcal A\xrightarrow{d_{\PP^4}(F_0+tG)}\mathcal B).
\]
Thus $\mathcal K=T_{\mathcal X/\Delta}$, and set
$E_0=\mathcal K|_{X_0}$. Write $Q=\bigoplus_{i=1}^{16}\CC_{p_i}$.

\begin{proposition}\label{prop:tangent-sheaf}
The sheaf $\mathcal K$ is flat over the disc and restricts to $T_{X_t}$
for $t\ne0$. There are identities and an exact sequence
\begin{equation}\label{eq:tangent-specialization}
 0\to E_0\to T_{X_0}\to Q\to0,\qquad
 E_0^{**}=T_{X_0},\qquad q_*T_Y=T_{X_0},\qquad
 R^jq_*T_Y=0\quad(j>0).
\end{equation}
At each node $T_{X_0}$ is reflexive of depth two, whereas $E_0$ has
depth one and is not reflexive. On each exceptional curve,
\begin{equation}\label{eq:tangent-splitting}
 T_Y|_{C_i}=\OO(2)\oplus\OO(-1)^{\oplus2}.
\end{equation}
\end{proposition}
\begin{proof}
Let $\mathcal I$ and $\mathcal Q$ be the image and cokernel of the
relative Jacobian map. Locally the total space is smooth and the
projection has equation $\sum z_j^2=t$. Hence $\mathcal Q$ is a
length-one skyscraper at each node, annihilated by $t$, and vanishes
elsewhere near the central fibre. Both $\mathcal I\subset\mathcal B$
and $\mathcal K\subset\mathcal A$ have no base-parameter torsion,
so they are flat over the smooth one-dimensional base. Tensor the
sequences
\[
 0\to\mathcal K\to\mathcal A\to\mathcal I\to0,
 \qquad 0\to\mathcal I\to\mathcal B\to\mathcal Q\to0
\]
with the central fibre. The first stays left exact, and the second
gives $\ker(\mathcal I_0\to\mathcal B_0)=
\operatorname{Tor}_1^{\OO_\Delta}(\mathcal Q,\CC)=Q$.
The kernel of $\mathcal A_0\to\mathcal B_0$ is $T_{X_0}$, so the
snake lemma gives the first sequence in \eqref{eq:tangent-specialization}.

At a node let $R=\CC\{z_1,\ldots,z_4\}/(\sum z_j^2)$ and let
$\mathfrak m$ be its maximal ideal. The tangent module is the kernel
of $R^4\to\mathfrak m$. The sequence
$0\to\mathfrak m\to R\to\CC\to0$ gives depth one for $\mathfrak m$,
and the kernel consequently has depth two. It is torsion-free and
locally free off the node, hence reflexive on this normal threefold.
The sequence with nonzero length-one quotient then gives depth one
for $E_0$. Both sheaves agree off codimension three, so their reflexive
hulls agree.

The exceptional locus in the smooth $Y$ has complex codimension two. Holomorphic sections of its tangent bundle extend
across those curves. On the common regular locus this identifies
$q_*T_Y$ with the reflexive extension $T_{X_0}$
\cite[Tags 0EBJ and 0AY6]{stacks}.
For the higher direct images, the tangent-normal sequence over $C_i$
splits because $H^1(\PP^1,\OO(3))=0$, giving
\eqref{eq:tangent-splitting}. If $I$ is the exceptional ideal,
$I^n/I^{n+1}=\OO(n)^{\oplus(n+1)}$. Tensoring with the tangent
restriction gives degrees $n+2$ and $n-1$, whose first cohomology
vanishes for every $n\ge0$. Induction through the thickening sequences
therefore gives $H^1(T_Y/I^{n+1}T_Y)=0$.
In an incidence chart the node coordinates are the entries of
$\left(\begin{smallmatrix}a&b\\ua&ub\end{smallmatrix}\right)$;
the pulled-back node ideal is $(a,b)=I$. These are exactly the
thickenings required by formal functions, not merely an equivalent
informal neighborhood model. Formal functions and faithful flatness
of completion give $R^1q_*T_Y=0$ \cite[Tag 02OC]{stacks}.
The fibres have dimension at most one, so all higher direct images
vanish as well.
\end{proof}

The exceptional splitting rules out locally free descent across the
contracted point. In particular the resolution direct image is the
reflexive hull, not the flat central fibre $E_0$. Compatibility is
through the common regular tangent bundle and the specified coherent
extensions. The exterior identifications on the smoothing side are
differentiable, not generally biholomorphic. No single flat family
containing both the resolution and smoothing branches is asserted.

Every smooth deformation has its intrinsic relative tangent bundle.
Thus this bundle choice introduces no additional existence obstruction
on smooth base deformations. The singular flatness calculation above
uses the actual one-dimensional family with nonzero linear local
smoothing terms. It is not a flatness theorem over an arbitrary
higher-dimensional singular deformation base. As a consistency check,
proper flatness, \eqref{eq:tangent-specialization} and the direct images
imply the holomorphic Euler-characteristic identity
$\chi(Y,T_Y)-\chi(X_t,T_{X_t})=16$, distinct from the topological
Euler-characteristic difference $32$.

\subsection{Actual metrics and controlled local curvature}
On the resolution use the K\"ahler class $H+\epsilon^2L$, the restriction
of a positive product class, and on the smoothing use $H$ with
$t=\epsilon^3$. Denote either branch by $Z_\epsilon$, with its actual
Ricci-flat metric $g_\epsilon$.

\begin{proposition}\label{prop:tangent-metrics}
The nontrivial tangent bundles $T_Y$ and $T_{X_t}$ are stable for these
classes. Their Ricci-flat tangent connections are irreducible zero-slope
Hermitian--Einstein connections. Both branches converge to the same
metric completion $(X_0,d_{g_0})$, homeomorphic to $X_0$, smoothly on
fixed regular compact subsets. Near each node, after rescaling by
$\epsilon^{-2}$, the metrics and tangent connections converge with all
derivatives to the normalized nonflat resolved-conifold or Stenzel
model.

For each node there is a collar radius $r$, extended over the cap with
$r+\epsilon\asymp\epsilon$ there, and fixed constants such that
\begin{equation}\label{eq:tangent-bounds}
 \begin{aligned}
 |\nabla^j\Rm(g_\epsilon)|&\le C_j(r+\epsilon)^{-2-j},\qquad j\ge0,\\
 \vol\{r<\rho\}&\le C\rho^6,\qquad
 \int_{\{r<\rho\}}|F_{T_{Z_\epsilon}}|^2dV\le C\rho^2,
 \qquad \epsilon\le\rho\le\rho_0,\\
 d_{GH}((Z_\epsilon,g_\epsilon),(X_0,d_{g_0}))&\le C\epsilon^\theta
 \quad\text{for some }\theta>0.
 \end{aligned}
\end{equation}
All tensor norms, derivatives and volumes in this display use the actual
metric. In particular the unscaled tangent curvature-energy measure has
no atom at a node.
\end{proposition}
\begin{proof}
Existence of the smooth Ricci-flat metrics is Yau's theorem \cite{yau}.
The holonomy argument following \eqref{eq:central-cohomology} gives
full $\mathrm{SU}(3)$. Ricci-flatness makes the tangent connection
Hermitian--Einstein of zero slope; irreducibility and the
Hermitian--Einstein stability implication give stability
\cite[Theorem~5.8.3]{kobayashi}. These bundles are nontrivial since their
second Chern pairing with $H$ is $50$.

The actual smoothing expansion was verified in
Proposition~\ref{prop:metric-comparison}. On the resolution all local
areas are $\int_{C_i}(H+\epsilon^2L)=\epsilon^2$. Thus the local AC
K\"ahler classes lie in the restriction image of the single global
class $L$, exactly Assumption R of Benabida's Theorem~B
\cite{benabida}; his Example~1.5 also treats this split quintic.
The face restrictions give the claimed central and rescaled local
limits, with all derivatives. For tangent connections these follow
by differentiation of the actual metrics and complex structures.

The positive face metrics and the bounded differentiated coefficients
in conic frames give uniform cone comparison on annuli and scaled
AC comparison on the core. Every physical derivative costs
$(r+\epsilon)^{-1}$; two derivatives give curvature. This proves the
first line of \eqref{eq:tangent-bounds} and core volume $O(\epsilon^6)$.
Integrating the annular volume comparison gives the displayed
$C\rho^6$ bound. On a dyadic annulus of radius $s\ge\epsilon$,
the tangent curvature energy is bounded by $Cs^{-4}s^6=Cs^2$.
Summing to the core, whose energy is $O(\epsilon^2)$, gives $C\rho^2$.
The central cone comparison adds one point in the metric completion
at each node and identifies its topology with the algebraic cone,
so the completion is homeomorphic to $X_0$.

We give the global distance estimate explicitly. Each filled region
$\{r<\rho\}$ has intrinsic diameter at most $C\rho$: radial annular
lengths sum to $O(\rho)$, fixed-radius links have diameter $O(\rho)$,
and the core has diameter $O(\epsilon)$. Collapsing the sixteen
regions to points changes Gromov--Hausdorff distance by at most
$C\rho$. Indeed a shortest quotient chain need visit each collapsed
point at most once, and each jump can be replaced by a path of its
intrinsic diameter. The same applies to the central space. On the
common exterior the actual tensor comparison, and its resolution
analogue from Theorem~B, give relative error $O((\epsilon/\rho)^{a_0})$.
The resulting quotient spaces have uniformly bounded diameter and
are bilipschitz with factor $1+O((\epsilon/\rho)^{a_0})$.
Consequently
\[
 d_{GH}(Z_\epsilon,X_0)\le C\bigl(\rho+(\epsilon/\rho)^{a_0}\bigr).
\]
Taking $\rho=\epsilon^{a_0/(1+a_0)}$ proves the last line of
\eqref{eq:tangent-bounds}, with a possibly smaller common exponent
for the two branches.
\end{proof}

The energy statement allows nonflat rescaled limits and unbounded
pointwise curvature. It supplies neither uniform estimates for arbitrary
gauge bundles nor a theorem for their HYM flows. Its role here is to
identify the actual tangent geometry compatible with the explicitly
calculated transition.

\subsection{Proof of the integrated theorem}
\begin{proof}[Proof of Theorem~\ref{thm:integrated}]
Proposition~\ref{prop:split-existence} establishes the genuine projective
family, the chosen resolution and the smooth fibres. The complete
intersection calculation in Section~\ref{sec:quintic} gives the
torsion-free smooth rings, their Chern classes and pairings. The Thom
and Mayer--Vietoris calculations there determine the exterior ring and
all cohomological attachments. The contraction and Milnor collapse
give the central ring and specialization maps by the relative pair
sequences, with the naturality described in
Section~\ref{sec:specialization}. This proves part~(1).

Use the collar-compatible trivialization of
Proposition~\ref{prop:metric-comparison} as the exterior identification
in the attachment calculation. Its horizontal lift agrees with the
prescribed node collars and is defined over the fixed ray. Two such
extensions away from the collars are isotopic along that ray, so they
induce the same marked cohomology maps. Thus the analytic and
topological link generators can be chosen identically.
Proposition~\ref{prop:restriction-spectrum} gives
\eqref{eq:split-exact-sequence}, since the hyperplane class $H$ is
primitive on the smooth quintic and $r(L)=\mathbf1$.
The freeness of $H^3(U;\ZZ)$ makes its vanishing sublattice primitive.
The disjoint vanishing spheres make it isotropic. The calculated
factorization $\Lambda_0\hookrightarrow\Lambda\twoheadrightarrow W$ gives
the displayed Picard--Lefschetz block and its Smith form, proving
part~(2).

For part~(3), the same integral class $L$ has
$L\cdot C_i=1$ and cubic data $(H^3,LH^2,L^2H)=(5,4,0)$.
It determines the exceptional areas $s^2$ in class $H+s^2L$.
On the central regular locus, Proposition~\ref{prop:mass} represents
$L-4H/5$ by the primitive harmonic form of squared norm $16/5$.
Identity~\eqref{eq:period-capacity-identification} identifies its
sixteen boundary coefficients with the same $r(L)=\mathbf1$.
The spectral theorem therefore gives $10\pi$, with its stated
count, type and error. Proposition~\ref{prop:tangent-metrics}
supplies the actual metric limits on the two branches. Apply its
resolution assertion with $\epsilon=s$ and its smoothing assertion
with $\epsilon=t^{1/3}$ separately; no ratio of $s$ and $t$ is used.

Finally, Proposition~\ref{prop:tangent-sheaf} identifies the flat
central tangent sheaf and its distinct reflexive hull, and
Proposition~\ref{prop:tangent-metrics} proves the tangent stability
and actual connection limits. This proves part~(4). These statements
use the same central regular tangent bundle; they do not identify
the flat central sheaf with the resolution direct image, or assert a
flat family containing both branches. The restriction, primitive norm
and actual collar identifications above provide the compatibility
between the topological and analytic assertions.
\end{proof}
